% bigPicture.tex — the overview document (read BEFORE the workbook)
%
% USAGE BY CLAUDE:
%   - READ THIS FIRST, before brief.tex and workbook.tex, at the start of a
%     session on a large project. It is the five-minute on-ramp: the goal, the
%     one organising idea, what follows from it, and an honest status ledger.
%   - It is equation-LIGHT on purpose. No derivations live here. Every claim
%     points to workbook.tex \S (section) or eq. for the details.
%   - The three-tier reading order for a big project is:
%         bigPicture.tex   (this file — the shape, in 5 minutes)
%       → brief.tex        (the current results, stated cleanly)
%       → workbook.tex     (the full derivations, navigated by grep)
%
% KEEPING IT HONEST:
%   - When a major result lands, update the status ledger below IN THE SAME
%     COMMIT as the workbook change. A stale ledger is worse than none.
%   - Quote workbook.tex's REAL section titles in the pointers. A plausible-but-
%     invented section name is worse than no pointer (a common failure when a
%     weaker model drafts the overview — check the pointer resolves).
%   - Do not overstate status. If Layer 1 proves only property P, say "P", not
%     "the full theorem". The ledger's whole job is to stop that inflation.
%
\documentclass[11pt, a4paper]{article}

\usepackage{amsmath, amssymb}
\usepackage{mathtools}
\usepackage[margin=2.5cm]{geometry}
\usepackage{microtype}
\usepackage[T1]{fontenc}
\usepackage[utf8]{inputenc}
\usepackage[colorlinks=true, linkcolor=blue, citecolor=blue, urlcolor=blue]{hyperref}
\usepackage{enumitem}
\usepackage{xcolor}
\usepackage[framemethod=default]{mdframed}

% -----------------------------------------------------------------------
% STATUS-LEDGER BOXES
% green  = proven (theorem / machine-checked)
% orange = open / numerically known but not proven
% Change the frame titles freely; keep the two colours distinct so status is
% readable at a glance.
% -----------------------------------------------------------------------
\newmdenv[
  linecolor={green!45!black}, backgroundcolor={green!6},
  linewidth=0.8pt, roundcorner=3pt,
  frametitle={\bfseries\color{green!45!black} Proven (theorems, machine-checked)},
  frametitlerule=true, frametitlerulecolor={green!45!black},
  frametitlebackgroundcolor={green!12},
  skipabove=10pt, skipbelow=10pt
]{provenbox}

\newmdenv[
  linecolor={orange!70!black}, backgroundcolor={orange!6},
  linewidth=0.8pt, roundcorner=3pt,
  frametitle={\bfseries\color{orange!70!black} Open / numerically known (not yet proven)},
  frametitlerule=true, frametitlerulecolor={orange!70!black},
  frametitlebackgroundcolor={orange!12},
  skipabove=10pt, skipbelow=10pt
]{openbox}

% -----------------------------------------------------------------------
\title{[Project Title] — Big Picture\\
{\normalsize\textit{The shape of the project in five minutes — read before the workbook}}}
\author{[Author Name]}
\date{\today}
% -----------------------------------------------------------------------

\begin{document}
\maketitle

% -----------------------------------------------------------------------
\section*{What we are trying to do}
% -----------------------------------------------------------------------
% One paragraph, plain language. The goal / the open question. No jargon that
% is not defined in the workbook's crash-course appendices.

[State the goal in one paragraph: what object, what question, what would count
as success. Write it so a colleague in a neighbouring field understands the aim
without reading anything else.]

% -----------------------------------------------------------------------
\section*{The one organising idea}
% -----------------------------------------------------------------------
% The single structure the whole project hangs on. If a reader remembers only
% one thing, this is it.

[Name the central structure / method / decomposition that everything else is
built on. One or two paragraphs. Point to where it is set up: workbook.tex
\S``[real section title]''.]

% -----------------------------------------------------------------------
\section*{What follows from it}
% -----------------------------------------------------------------------
% The consequences, as a short narrative — how the pieces connect toward the
% goal. Still equation-light; each step names its workbook location.

[Sketch the route from the organising idea to the goal. Two or three short
paragraphs. For each claimed step, cite workbook.tex \S/eq. rather than
reproducing it here.]

% -----------------------------------------------------------------------
\section*{Where we stand — status ledger}
% -----------------------------------------------------------------------
% The single glanceable source of "where are we". Keep it in sync with the
% workbook. Every line names its real workbook location.

\begin{provenbox}
\begin{itemize}[leftmargin=1.4em, itemsep=3pt]
  \item [Result that is fully proven] \quad(workbook.tex \S``[real section title]'').
  \item [Another proven result] \quad(workbook.tex \S``[...]'').
\end{itemize}
\end{provenbox}

\begin{openbox}
\begin{itemize}[leftmargin=1.4em, itemsep=3pt]
  \item [Result verified numerically to N digits — structure clear, value not yet
        analytic] \quad(workbook.tex \S``[...]'').
  \item [What is still missing for full rigour, stated precisely] \quad(needs [X]).
\end{itemize}
\end{openbox}

% -----------------------------------------------------------------------
\section*{Where the details live}
% -----------------------------------------------------------------------

\begin{itemize}[leftmargin=1.4em, itemsep=2pt]
  \item \textbf{Current results, stated cleanly:} \texttt{brief.tex}.
  \item \textbf{Full derivations and failed attempts:} \texttt{workbook.tex}
        (navigate by \texttt{grep}/section label; do not read top to bottom).
  \item \textbf{Definitions of the machinery / jargon:} the crash-course
        appendices in \texttt{workbook.tex}.
\end{itemize}

\end{document}
